\documentclass[10pt,a4paper]{../_shared/altacv}

\geometry{left=1.25cm,right=1.25cm,top=0.5cm,bottom=1cm,columnsep=1.2cm}
\usepackage[utf8]{inputenc}
\usepackage[T1]{fontenc}
\usepackage[french]{babel}
\usepackage{paracol}
\usepackage{xcolor}
\usepackage{hyperref}
\usepackage{fontawesome5}
\usepackage{setspace}

% Couleurs exactes du CV
\definecolor{SlateGrey}{HTML}{2E2E2E}
\definecolor{LightGrey}{HTML}{666666}
\definecolor{DarkPastelRed}{HTML}{8AC0D9}
\definecolor{PastelRed}{HTML}{00B0DA}
\definecolor{GoldenEarth}{HTML}{B4AD0C}
\colorlet{name}{black}
\colorlet{tagline}{PastelRed}
\colorlet{heading}{DarkPastelRed}
\colorlet{headingrule}{GoldenEarth}
\colorlet{subheading}{PastelRed}
\colorlet{accent}{PastelRed}
\colorlet{emphasis}{SlateGrey}
\colorlet{body}{LightGrey}

\renewcommand{\familydefault}{\sfdefault}

% En-tête strictement identique
\name{BOILEAU Guillaume}
\tagline{R\&D Engineer and Data Project Manager}
\photoL{2.8cm}{../_shared/moi} 
\personalinfo{
  \email{guillaume.boileaupro@gmail.com}
  \phone{+33 6 59 94 85 84}
  \location{Alpes Maritimes, France}
  \linkedin{guillaume-boileau-phd-891b81265}
  \researchgate{Guillaume-Boileau-2}
  \orcid{0000-0002-3576-6968}
}

% Espacement et lisibilité
\setlength{\parskip}{0.75\baselineskip}
\linespread{1.15}

\begin{document}

\makecvheader

% Active le grand trait vertical doré sans rien casser visuellement
\cvsection{Lettre de motivation}

\begin{paracol}{1}

%{\color{accent}\textbf{Objet :}} \textit{Candidature au poste d'ingénieur développement,  }

%\vspace{0.5cm}

{\color{body}
\iffalse
Nice, le 17 juillet 2025

Madame, Monsieur,

Titulaire d'un doctorat en physique et fort de plus de sept années d'expérience en ingénierie des données et en modélisation scientifique dans des environnements complexes, je souhaite aujourd'hui mettre mon expertise au service d'une structure dynamique et ambitieuse telle que Capgemini Engineering. Votre offre pour un poste d'Ingénieur DevOps m'a particulièrement intéressé, tant par sa dimension technique que par sa transversalité entre R\&D, automatisation et fiabilité opérationnelle.

Mon expérience dans des projets scientifiques complexes, notamment au sein du CNES et de la mission LISA (ESA/NASA), m'a permis de développer une expertise poussée dans la conception, le déploiement et la supervision d'architectures de traitement de données scientifiques à large échelle. J'ai notamment contribué à la mise en œuvre de pipelines d'analyse complexes intégrant Python, Bash, Git, Slurm, Condor et divers outils de supervision. Travaillant en lien étroit avec des ingénieurs système, des experts en sécurité et des équipes de développement, j'ai appris à concevoir des solutions robustes, automatisées et documentées dans des environnements distribués critiques.

Je possède de solides compétences sur les systèmes Linux, les protocoles réseaux (SSH, VPN, transfert sécurisé), ainsi qu'une bonne compréhension des logiques de CI/CD, de gestion de configuration et de virtualisation. J'ai également une appétence naturelle pour l'automatisation et la documentation technique, ainsi qu'une capacité à interagir efficacement avec les équipes techniques et non techniques.

Intégré au sein de grandes collaborations internationales (LIGO/Virgo/KAGRA), j'ai acquis une capacité à travailler dans des contextes exigeants, avec des enjeux de qualité, de performance et de disponibilité proches de ceux rencontrés dans les environnements DevOps avancés. Le poste proposé par Capgemini Engineering représente pour moi une évolution naturelle, me permettant de mettre à profit mes compétences en architecture système, en automatisation et en coordination technique.

Basé à proximité immédiate de Valbonne, je suis disponible à partir de septembre ou octobre. Votre approche centrée sur l'expertise, l'accompagnement personnalisé et l'engagement collectif me motive particulièrement. Je serais très heureux de pouvoir contribuer activement à vos projets.

Je reste bien entendu à votre disposition pour tout échange complémentaire et vous remercie pour l'attention portée à ma candidature.

Veuillez agréer, Madame, Monsieur, l'expression de mes salutations distinguées.

\textbf{Guillaume Boileau}

\fi
\textbf{Objet : Candidature au poste de Scientifique spécialiste de données radioastronomiques (H/F)} \\


Madame, Monsieur,

Titulaire d’un doctorat en physique, je possède plus de sept années d’expérience en ingénierie des données scientifiques, en modélisation statistique et en traitement du signal, acquises au sein de collaborations internationales de grande envergure. Je souhaite aujourd’hui mettre mon expertise au service du projet SPECTRUM, dont l’ambition de construire un continuum européen des données et du calcul scientifique résonne pleinement avec mon parcours.

Au CNES, j’ai conçu une plateforme de simulation modulaire en Python pour la mission spatiale LISA (ESA/NASA), intégrant des modèles complexes dans une chaîne de traitement robuste et pérenne. Ce pipeline, validé comme socle pour le segment sol de la mission, illustre ma capacité à structurer des outils techniques au cœur d’une architecture scientifique distribuée. J’ai travaillé en lien direct avec les équipes d’ingénierie système et les experts métier, en assurant la coordination technique et la documentation des livrables critiques.

J’ai également contribué au développement de chaînes bayésiennes (MCMC) pour l’atténuation des bruits instrumentaux à l’Université d’Anvers, dans un contexte de capteurs haute-précision. Cette expérience m’a permis de maîtriser des infrastructures de calcul distribuées (\texttt{Condor}, \texttt{Slurm}), l’automatisation sous \texttt{Linux}, et les outils de développement collaboratif (\texttt{Git}, intégration continue). J’accorde une grande importance à la reproductibilité, à la supervision système et à la qualité du code dans les environnements scientifiques complexes.

Mon implication active dans les collaborations LISA, LIGO/Virgo/KAGRA et Einstein Telescope m’a familiarisé avec la coordination de communautés scientifiques internationales, la gestion de projets transverses et la co-construction de livrables stratégiques. J’ai encadré des stagiaires, dirigé des équipes techniques, et participé à l’élaboration de solutions logicielles pour des plateformes de traitement à l’échelle.

Je souhaite mettre mon expérience au service du projet SPECTRUM en contribuant à l’animation de la communauté radioastronomique, à la collecte structurée des cas d’usage et des exigences, et à la consolidation des éléments techniques et stratégiques du SRIDA. Travaillant déjà à l’UMR Lagrange, je suis particulièrement sensible aux enjeux de fédération des données et d’harmonisation des pratiques entre disciplines, et motivé à y jouer un rôle actif.

Je vous remercie pour l’attention portée à ma candidature, et me tiens à votre disposition pour tout échange complémentaire.

\vspace{0.5cm}

Veuillez agréer, Madame, Monsieur, l’expression de mes salutations distinguées.

\vspace{0.5cm}

Guillaume Boileau

}


\end{paracol}

\end{document}
